
\section{Introduction}\label{sec:intro}

A causal discovery algorithm returns a graph, and a practitioner reading it
wants to know which of its features the data support. Constraint-based methods
such as PC and FCI
\citep{spirtes2000causation,spirtes1995causal,colombo2012learning} start from
the complete graph and delete an edge whenever a conditional-independence test
fails to reject. Failure to reject is compatible with independence, with a
dependence too weak to see at the current sample size, and with a test that had
little power; the deletion happens regardless. Two consequences follow.
Finite-sample correctness requires faithfulness or a quantitative strengthening
of it \citep{robins2003uniform,uhler2013geometry}, because only such an
assumption turns ``the test did not reject'' into ``the variables are
independent''. And the output is two-valued where the evidence is three-valued:
a missing edge means either ``the data indicate no edge'' or ``nothing was
established about this pair''.

There is a third consequence that matters in practice. Data accumulate, and an
analyst who inspects the estimate as they do, and keeps any edge that has ever
appeared, is performing optional stopping. A fixed-sample guarantee does not
cover this. Section~\ref{sec:exp-monitoring} measures the effect on eight-node
graphs: when PC, GES and DirectLiNGAM are rerun at each sample size from $100$
to $6400$, between a quarter and nine tenths of the graphs have at least one
false edge in the union of the estimates.

This paper studies the opposite discipline, which we call \emph{certificate-only}
discovery: assert a feature of the graph only by rejecting a null hypothesis
that is true whenever the feature is absent, and never assert anything by
failing to reject. An algorithm that obeys the rule starts from the empty graph
and adds; it cannot report an edge as absent, only as undecided; and what it
controls is the probability of asserting something false. If each test is an
anytime-valid process, the guarantee holds at every sample size at once and
survives continuous monitoring.

\paragraph{Contributions.}

\begin{enumerate}
\item \emph{Adjacency certificates under monitoring.} For a linear structural
  equation model with independent errors, d-separation implies that a partial
  regression coefficient is zero, so adjacency has a certifying null that
  needs no faithfulness and no Gaussian errors. The minimum over conditioning
  sets of per-set sequential tests bounds the probability that a false edge is
  \emph{ever} certified (Section~\ref{sec:adjacency}). The bound is exact for
  Gaussian errors and asymptotic otherwise, and we measure how far from
  nominal it is off-model (Section~\ref{sec:exp-primitive}).
\item \emph{An orientation certificate based on predictive fit.} For a pair of
  variables we fit forecasters for the factorisations $x\to y$ and $y\to x$
  sequentially, and certify a direction when the running mean of the
  log-score difference has a confidence sequence that excludes zero under two
  model classes at once (Section~\ref{sec:orientation}). The claim is about
  predictive fit and is valid whether or not either model is correct; it
  identifies the causal direction under additive-noise models.
\item \emph{Evidence under pre-specified protocols.} We compare against PC,
  GES and DirectLiNGAM under monitoring (Section~\ref{sec:exp-monitoring}),
  evaluate the orientation certificate on simulated pairs and on the
  Tübingen cause-effect pairs with a protocol committed before the run
  (Sections~\ref{sec:exp-pairs} and~\ref{sec:exp-tuebingen}), evaluate it on physical devices with a known graph
  (Section~\ref{sec:exp-chamber}), study scaling with graph size, sample size and the constants of the algorithm
  (Sections~\ref{sec:exp-scaling} and~\ref{sec:exp-ablation}), and report the
  regimes where the certificates fail (Sections~\ref{sec:exp-pairs}
  and~\ref{sec:exp-scope}).
\end{enumerate}

\paragraph{Scope.} The method targets continuous, observational data from a
causally sufficient system whose mechanisms are linear or additive-noise, at
sample sizes of a few hundred or more. We chose the benchmarks to match this
domain and fixed the selection criteria before running; data outside it are
reported as a limitation, not dropped.

\paragraph{Outline.} Section~\ref{sec:related} places the work in the literature.
Section~\ref{sec:method} develops the certificates and the algorithm.
Section~\ref{sec:experiments} reports the experiments and
Section~\ref{sec:discussion} discusses limitations. Proofs and further results
are in the appendices.
